\section{Introduction}\label{sec:intro}

Applied causal work now spans a large and fragmented software stack:
\proglang{Stata} commands, methodological-reference \proglang{R}
packages, and an expanding but incomplete \proglang{Python} ecosystem.
A single project can require OLS and IV, staggered
difference-in-differences, regression discontinuity, synthetic control,
matching, double machine learning, sensitivity analysis, and publication
tables. The statistical literature has converged toward more explicit
design checks and estimator-specific inference
\citep{imbens2009recent,angrist2009mostly,callaway2021difference,
roth2023whats,wing2024designing,baker2026difference,rambachan2023more,abadie2010synthetic,
chernozhukov2018double}; the software interface has not converged with
it. Users still cross package boundaries, result classes, naming
conventions, and diagnostic vocabularies. Large-language-model agents
add a further constraint: they need typed schemas and structured errors,
not prose-only help pages.

This paper introduces \statspai{}, a \proglang{Python} platform for
applied econometrics and causal inference. The package provides
(i) a unified \code{import statspai as sp} surface across common causal
designs, (ii) shared reporting conventions for mature estimator results,
diagnostics, citations, and export methods, (iii) a machine-readable registry that can be exposed through a
Model Context Protocol (MCP) server~\citep{anthropic2024mcp}, and (iv) an explicit validation
discipline against analytical targets, canonical \proglang{R}
implementations, and selected \proglang{Stata} migration bridges. This article describes \statspai{} \StatsPAIVersion{}, distributed on PyPI
(\code{pip install statspai}, tag \code{v\StatsPAIVersion{}}) under the MIT
licence, with source at
\url{https://github.com/brycewang-stanford/StatsPAI}.

A short paper on the package appeared in the \emph{Journal of Open
Source Software}~\citep{wang2026statspaijoss}. It established the
unified interface, the registry and agent schemas, the structured result
objects, and the existence of a same-byte cross-language parity harness
and of a 1{,}000-replication coverage run for OLS, DiD, and IV. This
article does not re-establish those facts. Its question is the one a
parity harness raises but cannot answer by itself: \emph{what exactly
does a validation label certify?} A green parity row can be earned by
the wrong code (a \proglang{Python} wrapper that calls the \proglang{R}
reference it is compared against), vouch for the wrong function (an
alias credited with evidence that ran another entry point), rest on the
wrong statistical comparison (a stochastic estimator judged against its
sampling standard error rather than its own Monte Carlo error), or
describe a configuration the user is not running. The contribution is
the design that makes those failure modes checkable, and the evidence
of what applying it found. Concretely, the article adds (i) evidence
provenance that is verified rather than asserted --- every parity row is
classified by what the \statspai{} side actually executes, checked by
a call trace, and every inherited label must name a test that ran the
inheriting function (Sections~\ref{sec:track-a}
and~\ref{sec:attachment-audit}); (ii) statistical-applicability
semantics for automated diagnostics, so that a check which does not
exist for a fit is reported as not applicable rather than missing
(Section~\ref{sec:ex-card}); (iii) a seed-replicated standard for
stochastic estimators that separates algorithmic Monte Carlo error from
sampling error and states equivalence against an explicit margin
(Section~\ref{sec:validation-tiers}); (iv) a fixed adjudication procedure
for disagreements, with its outcomes (Section~\ref{sec:adjudication});
and (v) new evidence the JOSS paper does not contain: the estimand and
identification contracts of a twelve-estimator core, configuration-level
coverage diagnostics (bias, Monte Carlo SD, standard-error calibration,
interval length), a performance study that reports unfavourable
comparisons, and seven executed worked examples. Applying the design to
its own package found defects that numerical parity alone had not: two
Honest-DiD rows that compared \proglang{R} with itself, a causal-forest
coverage row whose entry point reported a separate AIPW estimator rather
than the forest's, an entropy-balancing standard error twice the
estimator's Monte Carlo SD, and an audit that asked for an
over-identification test on a just-identified model.
The registry distinguishes \code{certified}, \code{validated},
\code{api\_stable}, and \code{experimental} symbols -- queryable at
call time -- so numerical validation is an
evidence tier rather than a blanket claim over the whole API surface;
the full census opens Section~\ref{sec:parity}.

Listing~\ref{lst:motivating} fixes ideas. The classical
\citet{card1995using} returns-to-schooling design is estimated four
ways --- OLS, instrumental variables, double machine learning
\citep{chernozhukov2018double}, and a causal forest
\citep{wager2018estimation} --- from a single import, and three of the
four fitted objects, which belong to three different result classes,
are rendered into one publication table by a single call.
Three design properties recur throughout the paper: the
estimator surface is \emph{flat} (classical and modern methods are
siblings under \code{sp.}); the result objects are \emph{interoperable}
(mixed estimators feed directly into one \code{sp.regtable} call); and
the platform exposes \emph{audit hooks} (\code{sp.audit}) that list
the diagnostic evidence a design demands but the user has not yet
supplied. Section~\ref{sec:ex-card} runs the same session at a prompt
and prints what each of these calls returns.

Every listing in this paper is executed verbatim by the replication
archive's listing audit, and for the transcripts the audit goes
further: it replays each one statement by statement through a real
interpreter loop and compares the output printed here, line for line,
with the output the code actually produces
(Section~\ref{sec:reproducibility}). A printed transcript is quoted as
evidence and read as evidence, so it is checked like evidence.

\begin{lstlisting}[caption={\statspai{} reproduces and extends Card
(1995) from a single import; classical and modern estimators are
siblings under \code{sp.} and return one interoperable result
contract.},label={lst:motivating},captionpos=b]
import statspai as sp

card = sp.datasets.card_1995()
X    = ["exper", "expersq", "black", "south", "smsa"]

ols = sp.regress("lwage ~ educ + " + " + ".join(X), data=card)
iv  = sp.iv("lwage ~ " + " + ".join(X) + " + (educ ~ nearc4)",
            data=card)
dml = sp.dml(card, y="lwage", d="educ", X=X,
             model="plr", n_folds=5, random_state=42)
cf  = sp.causal_forest("lwage ~ educ | " + " + ".join(X), data=card,
                       n_estimators=2000, discrete_treatment=False,
                       random_state=42)

table = sp.regtable(ols, iv, dml, keep=["educ", "ATE"], stars=True,
                    model_labels=["OLS", "2SLS", "DML PLR"])
table.to_docx("returns_to_schooling.docx")
\end{lstlisting}

The package surface is deliberately broad, spanning classical
regression and panel models, the modern DiD, RD, and synthetic-control
literatures, ML-based treatment-effect estimation, matching and
weighting, sensitivity analysis, and publication output.
Breadth, however, is not the paper's evidentiary claim.
In this paper, ``validated'' means a
specific artifact exists: known-truth recovery, cross-language parity on
identical data bytes, Monte Carlo coverage, or an explicit convention
disclosure. Unit and regression tests are API-contract evidence; they
do not by themselves promote a helper into the numerical-validation
tier. An \code{api\_stable} helper has a stable call signature; it does
not acquire statistical credibility by proximity to validated estimators.
Table~\ref{tab:scope} makes this boundary explicit before the paper
turns to architecture and validation evidence.

\begin{table}[!ht]
\centering\scriptsize
\begin{tabular}{p{0.23\linewidth}p{0.36\linewidth}p{0.31\linewidth}}
\toprule
Layer & Examples & Evidence used in this paper \\
\midrule
Unified causal API & OLS/IV, DiD, RD, synth, DML, forests, matching,
causal-impact workflows & Shared entry points, mature result objects, exports,
and worked examples \\
Validated causal core & The twelve-estimator validation suite
& Analytical recovery, R parity, selected Stata bridges, coverage, and
performance checks \\
Agent-facing surface & Registry schemas, agent cards, MCP tools,
structured errors & Schema coverage, deterministic MCP trace, and
citation/audit checks \\
Broader registry & Long-tail econometric helpers and optional backends
& Reported as \code{api\_stable} or \code{experimental} unless backed by
the same evidence standard \\
\bottomrule
\end{tabular}
\caption{Scope of the submission. The package is broader than the
validated core; the paper-level claim is limited to the rows backed by
the artifacts in Sections~\ref{sec:examples}--\ref{sec:agentbench}.}
\label{tab:scope}
\end{table}

The closest implementations cover important parts of this terrain, but
with different centers of gravity; Table~\ref{tab:related-software}
states the comparison directly. Empirical comparisons appear later in
the same-byte parity rows
(Tables~\ref{tab:track-a-cross-language-snapshot}
and \ref{tab:headline-parity}) and in the timing study
(Figure~\ref{fig:track-c-loglog}); per-entry-point detail is generated
into \path{Paper-JSS/parity/replacement-table.md}. The contribution is therefore an
integration-and-validation layer around reference software, not a claim
that \statspai{} supersedes specialized packages or reference
implementations in their home domains, such as
\pkg{fixest}, \pkg{did}, \pkg{rdrobust}, \pkg{Synth},
\pkg{synthdid}, \pkg{DoubleML}, or \pkg{grf}. Representative software
citations include \pkg{Synth}, \pkg{MatchIt}, \pkg{DoubleML},
\pkg{EconML}, \pkg{CausalML}, and \pkg{DoWhy}
\citep{abadie2011synth,ho2011matchit,bach2022doubleml,econml,
chen2020causalml,sharma2020dowhy}.

\begin{table}[!ht]
\centering\scriptsize
\begin{tabular}{p{0.23\linewidth}p{0.32\linewidth}p{0.35\linewidth}}
\toprule
Comparable implementation & Where it remains the reference choice & \statspai{} contribution and boundary \\
\midrule
\pkg{statsmodels}, \pkg{linearmodels}, \pkg{pyfixest} &
Mature regression, panel, covariance, and HDFE APIs with focused
maintainer communities. &
\statspai{} adds a causal-workflow layer, audit prompts, citations, and
agent schemas around common calls; it is not a replacement for their
full model-specific APIs or fastest specialist paths. \\
\pkg{DoubleML}, \pkg{EconML}, \pkg{CausalML}, \pkg{grf} &
Deeper ML-causal and CATE-specific surfaces, tuning controls, and
method-family specialization. &
\statspai{} places DML and forests inside the same result, audit, and
validation ledger; the forest row is only a T3 stochastic agreement
claim, not deterministic parity with every \pkg{grf} option. \\
\pkg{DoWhy} &
Graph-first identification and refutation workflows. &
\statspai{} is estimator- and evidence-led: it emphasizes parity
artifacts and shared output contracts, but it does not attempt to match
graph-identification automation. \\
\pkg{did}, \pkg{fixest}, \pkg{rdrobust}, \pkg{rddensity},
\pkg{Synth}, \pkg{synthdid}, \pkg{MatchIt}, and Stata commands such
as \code{csdid}, \code{reghdfe}, \code{rdrobust}, and \code{sdid} &
Authoritative reference implementations in their home ecosystems,
often with more mature edge-case options and community history. &
\statspai{} validates against these references on identical data bytes
where possible and records disagreements as ledger entries rather than
strict parity claims; the remaining T4 disclosure is the classical-SCM
identification/non-uniqueness gap on the calibrated Basque replica,
while generalized SCM, RD-density,
and SDID are native parity rows in \StatsPAIVersion{}. Stata bridges require
a separate licensed runtime. \\
\pkg{differences}, \pkg{causalimpact}, and narrower workflow packages &
Focused workflows with less surface area for users who need exactly
that model family. &
\statspai{} standardizes cross-family handoffs through one namespace
and one post-estimation contract; that breadth comes with a larger
dependency surface and a validation ledger that should be read row by
row. \\
\bottomrule
\end{tabular}
\caption{Closest existing implementations and comparative scope,
surveyed in May 2026 and spot-checked in September 2026.}
\label{tab:related-software}
\end{table}

The submission should therefore be read in three layers. The manuscript
describes the user-facing API and the statistical-software design
choices; the validation tables state which methods have parity,
coverage, performance, or boundary evidence; and the replication archive
contains the executable ledgers that determine whether a label such as
\code{certified}, \code{validated}, or \code{api\_stable} is justified.
This separation is intentionally conservative: a broad namespace is
useful only if readers can tell which parts are statistical evidence and
which parts are stable software infrastructure.
To keep the printed article compact, the archive also includes a
reviewer evidence map and editor screening checklist that route
validation-scope, reproduction, related-review and COI, sustainability,
and release-boundary questions to executable artifacts.

Section~\ref{sec:architecture} describes the architecture the evidence
attaches to and Section~\ref{sec:agent} the agent-facing interface;
Section~\ref{sec:examples} runs seven worked examples, with staggered
DiD as the main case study; Sections~\ref{sec:parity}--\ref{sec:agentbench}
report the validation, performance, and interface evidence.
